\begin{ouverture}[calcul]{4}{Diviser et interpréter un partage}{Partie 1 · Nombres et calculs}
{Partage · Groupement · Reste · Division décimale}
\objectif{Distinguer le calcul d’une part et le calcul d’un nombre de parts.}
\objectif{Effectuer une division euclidienne et une division décimale posées.}
\objectif{Interpréter un quotient, un reste ou une valeur approchée.}
\prerequis{Tables de multiplication ; multiples et diviseurs ; division posée du CM2.}
\end{ouverture}
\pagestyle{otMain}\essentiel\fontsize{11}{15}\selectfont
\phantomsection\label{cours-4-1}
\essSection{1}{Deux sens pour une division}
\begin{definition}Une division répond à deux sortes de questions. Partage : 84 bonbons sont partagés également entre 6 enfants ; on cherche la valeur d’une part. Groupement : 84 bonbons sont rangés en sachets de 6 ; on cherche le nombre de parts. Le même calcul $84:6=14$ répond aux deux questions.\end{definition}
\label{fig-cours-4-1-cours-partage}\figureGeo{cours-partage}

\essSection{2}{Division euclidienne et reste}
\begin{definition}Effectuer la division euclidienne d’un entier $a$ par un entier $b$ non nul, c’est trouver le quotient $q$ et le reste $r$ tels que $a=b\times q+r$, avec $r<b$. Par exemple $487=23\times21+4$ et $4<23$ : le quotient est 21 et le reste est 4.\end{definition}
\begin{savoirfaire}[Encadrer par des multiples]Pour $487:23$, on encadre : $23\times20=460$ et $23\times30=690$, donc le quotient s’écrit avec deux chiffres et commence par 2. Puis $487-460=27$, qui contient encore une fois 23 : quotient 21, reste $27-23=4$.\end{savoirfaire}
\clearpage
\phantomsection\label{cours-4-2}
\essSection{3}{Division décimale par un entier}
\begin{savoirfaire}[Poursuivre après la virgule]Pour $23:4$ : 23 unités divisées par 4 donnent 5 unités ; il reste 3 unités, soit 30 dixièmes. 30 dixièmes divisés par 4 donnent 7 dixièmes ; il reste 2 dixièmes, soit 20 centièmes. 20 centièmes divisés par 4 donnent 5 centièmes, reste 0. Donc $23:4=5{,}75$.\end{savoirfaire}
\label{fig-cours-4-2-cours-potence}\figureGeo{cours-potence}

\begin{exemple}$47{,}6:4$ : on divise les unités, puis les dixièmes. On trouve $11{,}9$. Contrôle : $11{,}9\times4=47{,}6$.\end{exemple}
\essSection{4}{Résultat exact et division qui se poursuit}
\begin{propriete}Parfois la division s’arrête avec un reste nul : le quotient est exact. Parfois le même reste revient sans fin : $10:3=3{,}333\ldots$ On donne alors une valeur arrondie, par exemple $3{,}33$ au centième.\end{propriete}
\begin{attention}Le reste se lit dans l’histoire. 50 élèves partent en minibus de 12 places : $50=12\times4+2$. Il faut 5 minibus, car les 2 élèves restants doivent aussi voyager.\end{attention}

\clearpage\rubrique{Découverte}
\begin{blocLivre}{ex-4-1}\begin{exercice}[1]{Part ou nombre de parts ?}
Pour chaque problème, dire si l’on cherche la valeur d’une part ou le nombre de parts, puis le résoudre.\begin{questions}
\item 96 cartes sont partagées également entre 8 joueurs. Combien de cartes chacun reçoit-il ?
\item 96 cartes sont rangées en paquets de 8. Combien de paquets obtient-on ?
\end{questions}
\end{exercice}\end{blocLivre}
\begin{blocLivre}{ex-4-2}\begin{exercice}[1]{Division euclidienne}
Effectuer la division euclidienne de 487 par 23, puis de $1\,000$ par 7. Écrire chaque fois l’égalité obtenue.
\end{exercice}\end{blocLivre}
\begin{blocLivre}{ex-4-3}\begin{exercice}[1]{La bonne égalité}
Laquelle de ces égalités traduit la division euclidienne de 158 par 12 ? $158=12\times13+2$ ; $158=12\times12+14$ ; $158=12\times11+26$.
\end{exercice}\end{blocLivre}
\begin{blocLivre}{ex-4-4}\begin{exercice}[1]{Division décimale}
Poser et calculer $47{,}6:4$, puis $9{,}45:5$.
\end{exercice}\end{blocLivre}
\begin{blocLivre}{ex-4-5}\begin{exercice}[1]{Après la virgule}
Calculer le quotient décimal exact de $23:4$, puis de $7:8$.
\end{exercice}\end{blocLivre}

\rubrique{Entraînement}
\begin{blocLivre}{ex-4-6}\begin{exercice}[2]{Une division posée}
Recopier et compléter la division posée de 365 par 8, puis écrire l’égalité de la division euclidienne.
\label{fig-ex-4-6}\figureGeo{div-6}
\end{exercice}\end{blocLivre}
\begin{blocLivre}{ex-4-7}\begin{exercice}[2]{Le reste dans l’histoire}
50 élèves partent en sortie dans des minibus de 12 places. Combien de minibus faut-il ? Combien de places restent libres ?
\end{exercice}\end{blocLivre}
\begin{blocLivre}{ex-4-8}\begin{exercice}[2]{Une division sans fin}
Calculer $10:3$ et $2:3$ en poursuivant après la virgule. Que remarque-t-on ? Donner l’arrondi au centième de chaque quotient.
\end{exercice}\end{blocLivre}
\begin{blocLivre}{ex-4-9}\begin{exercice}[2]{Partager au centime}
25 € sont partagés également entre 4 amis, puis 17 € entre 6 amis. Les partages tombent-ils juste au centime près ?
\end{exercice}\end{blocLivre}
\begin{blocLivre}{ex-4-10}\begin{exercice}[2]{Retrouver le dividende}
Dans une division euclidienne, le diviseur est 15, le quotient 23 et le reste 9. Quel est le dividende ? Le reste aurait-il pu valoir 15 ?
\end{exercice}\end{blocLivre}

\rubrique{Problèmes}
\begin{blocLivre}{ex-4-11}\begin{exercice}[3]{Les boîtes d’œufs}
Une ferme ramasse $1\,000$ œufs dans la journée. Elle les range dans des boîtes de 6. Combien de boîtes pleines obtient-elle ? Combien de boîtes faut-il pour emballer tous les œufs ?
\end{exercice}\end{blocLivre}
\begin{blocLivre}{ex-4-12}\begin{exercice}[3]{L’addition du restaurant}
Sept amis partagent également une addition de $94{,}50$ €. Combien chacun paie-t-il ?
\end{exercice}\end{blocLivre}
\begin{blocLivre}{ex-4-13}\begin{exercice}[3]{La randonnée}
Une randonnée de $38{,}4$ km se fait en 3 jours de même longueur. Quelle distance par jour ? On la découpe ensuite en 8 étapes égales : quelle est la longueur d’une étape ?
\end{exercice}\end{blocLivre}
\begin{blocLivre}{ex-4-14}\begin{exercice}[3]{Quel jour sera-t-il ?}
Nous sommes lundi. Quel jour de la semaine serons-nous dans 100 jours ?
\end{exercice}\end{blocLivre}
\begin{blocLivre}{ex-4-15}\begin{exercice}[3]{Le plus grand reste}
Dans une division euclidienne par 24, quel est le plus grand reste possible ? Trouver le plus grand nombre inférieur à 500 qui donne ce reste.
\end{exercice}\end{blocLivre}

\clearpage
